\chapter{As entradas da máquina: como os olhos, os ouvidos e outros sensores estão ligados a ela}
\chaptermark{As entradas da máquina}
\label{sec:sensors}

\section{O sensor como transdutor}

Na fig.~\ref{fig:machine} os dados entravam na máquina pela esquerda, pelo bloco ``sensores''. Agora vamos abrir esse bloco. Para o engenheiro, cada órgão dos sentidos é um \emph{transdutor} com um estágio de entrada analógico e um conversor analógico-digital no fim; só que os dígitos na saída não são números, e sim disparos.

O esquema é o mesmo para todos os órgãos dos sentidos (fig.~\ref{fig:transducer}). Uma grandeza física (luz, pressão, vibração, uma molécula) atua sobre uma proteína receptora na membrana da célula sensorial. Essa proteína funciona como uma comporta: ela mesma, ou por uma curta cadeia de reações, abre um canal iônico. Surge uma \emph{corrente receptora} e, com ela, uma tensão analógica na membrana. Em seguida a tensão passa por um controle de ganho e chega ao codificador de disparos, o já conhecido neurônio integrador~\eqref{eq:glutneuron}, que transforma o nível em taxa e em instantes de disparo. A saída é quase sempre a mesma: os neurônios sensoriais do olho, do ouvido, do olfato e da pele liberam glutamato, de modo que as entradas se ligam diretamente ao barramento \nt{GLUT}.

\begin{figure}[t]
\centering
\begin{tikzpicture}[>=Stealth,
  blk/.style={draw,very thick,rounded corners=2pt,align=center,font=\footnotesize,minimum height=1.0cm,text width=1.9cm,fill=blue!5}]
\node[align=center,font=\small] (P) at (0.1,0) {luz, som,\\pressão,\\molécula};
\node[blk] (R) at (3.0,0) {receptor-\\comporta};
\node[blk] (G) at (6.5,0) {controle\\de ganho};
\node[blk] (E) at (10.0,0) {codificador\\de disparos};
\node[align=center,font=\small] (B) at (13.4,0) {barramento\\\nt{GLUT}};
\draw[->,very thick] (P) -- (R);
\foreach \ic/\xx in {sun/-1.1,speaker/-0.3,press/0.5,molecule/1.3}{\pic[scale=0.85] at (\xx,1.05) {\ic};}
\draw[->,very thick] (R) -- node[above,font=\scriptsize] {corrente} (G);
\draw[->,very thick] (G) -- node[above,font=\scriptsize] {nível} (E);
\draw[->,very thick,red!70!black] (E) -- node[above,font=\scriptsize,black] {disparos} (B);
\draw[decorate,decoration={brace,amplitude=5pt,mirror},gray!60!black] (1.8,-0.8) -- (7.7,-0.8) node[midway,below=5pt,font=\scriptsize,black] {parte analógica};
\draw[decorate,decoration={brace,amplitude=5pt,mirror},gray!60!black] (8.8,-0.8) -- (11.2,-0.8) node[midway,below=5pt,font=\scriptsize,black] {disparos};
\end{tikzpicture}
\caption{Qualquer órgão dos sentidos como uma cadeia de conversões. A proteína receptora abre um canal e gera corrente; o controle de ganho a adapta às circunstâncias; o codificador~\eqref{eq:glutneuron} transforma o nível em disparos, que seguem para o barramento \nt{GLUT}. No olho e no ouvido os primeiros estágios não emitem disparo algum: o sinal continua analógico até precisar ser levado longe.}
\label{fig:transducer}
\end{figure}

Um detalhe é bem conhecido de quem trabalha com eletrônica. No olho e no ouvido as primeiras células não emitem disparo algum: passam às vizinhas uma tensão contínua, como um estágio analógico numa placa. Os disparos só aparecem onde o sinal precisa ir longe. Um sinal analógico, ao longo de milímetros, se atenua e acumula ruído, enquanto um disparo, como vimos no capítulo~\ref{sec:neurontypes}, chega ao fim do fio com a mesma altura. Digitaliza-se onde começa a linha longa.

\section{No limite da física}

Os sensores do cérebro trabalham perto dos limites físicos de sensibilidade. O sensor de luz no escuro responde a \emph{um único fóton}: a absorção de uma só partícula de luz dá uma corrente de cerca de um picoampere, perceptível acima do seu próprio ruído~\cite{Baylor1979}. O sensor de som percebe um deslocamento do seu feixe sensível menor que um nanômetro~\cite{Hudspeth1989}.

Quando há pouca luz, a precisão não é limitada pelo sensor, e sim pela própria luz: os fótons chegam ao acaso, num fluxo de Poisson. Se durante o tempo de acumulação $T$ chegam ao sensor em média $N=\Phi T$ fótons, o seu número flutua em $\sqrt N$. Para que um acréscimo $\Delta N$ seja perceptível, ele deve ser algumas vezes maior que essa flutuação, e a fração distinguível de variação do brilho é
\begin{empheq}[box=\keybox]{equation}
\frac{\Delta I}{I} \;\approx\; \frac{k}{\sqrt{\Phi\,T}} .
\label{eq:photonnoise}
\end{empheq}
É a mesma fórmula da contagem de disparos~\eqref{eq:rateerror}: o ruído balístico (\emph{shot noise}) dos fótons e o dos disparos obedecem à mesma lei. No escuro, o olho só pode melhorar a precisão de um jeito: acumulando fótons por mais tempo, e o tempo de acumulação de fato cresce até alguns décimos de segundo. Por isso enxergamos mais devagar no crepúsculo.

\section{Faixa dinâmica: controle de ganho}

O problema de engenharia mais difícil dos sensores é a faixa. A iluminação, de uma noite estrelada até a neve ao sol, varia cerca de dez ordens de grandeza; a intensidade do som, do limiar de audição até o limiar da dor, doze. Já a taxa de disparos vai de zero a algumas centenas por segundo, menos de três ordens. Não dá para encaixar linearmente essa entrada nessa saída.

A solução é a mesma de qualquer receptor: \emph{controle automático de ganho}. As respostas das células sensoriais do olho são bem descritas pela fórmula
\begin{empheq}[box=\keybox]{equation}
r \;=\; r_{\max}\,\frac{I}{I+\sigma} ,
\label{eq:nakarushton}
\end{empheq}
em que $\sigma$ é o brilho em que a resposta chega à metade~\cite{NakaRushton1966}. Por si só,~\eqref{eq:nakarushton} cobre apenas duas ou três ordens. Mas $\sigma$ não é constante: com um trabalho prolongado sobre um fundo claro, ela cresce acompanhando o brilho médio. A curva de resposta se desloca ao longo do eixo logarítmico do brilho e sempre coloca o seu trecho íngreme onde a entrada está agora (fig.~\ref{fig:adaptation}). É a mesma divisão pela atividade total da inibição por derivação do capítulo~\ref{sec:gaba}~\cite{Carandini2012}, só que aqui o denominador é o brilho médio dos últimos segundos.

\begin{figure}[t]
\centering
\begin{tikzpicture}[>=Stealth,x=1.25cm,y=2.8cm]
\draw[->,gray!70!black] (-2,0) -- (6.4,0) node[below left,font=\small,black] {$\log_{10} I$};
\draw[->,gray!70!black] (-2,0) -- (-2,1.15) node[left,font=\small,black] {$r/r_{\max}$};
\foreach \u in {-2,0,2,4,6}{\draw[gray!60!black] (\u,-0.02) -- (\u,0.02); \node[below,font=\scriptsize] at (\u,0) {$\u$};}
\draw[densely dashed,gray!60!black] (-2,0.5) -- (6.2,0.5);
\foreach \s/\c/\lab in {0/blue!60!black/{fundo escuro},2/green!45!black/{médio},4/red!70!black/{claro}}{
  \pgfmathsetmacro{\lo}{max(-2,\s-3)}
  \draw[very thick,\c] (-2,0) -- (\lo,0) plot[domain=\lo:6.2,samples=80] (\x,{1/(1+10^(\s-\x))});
  \node[font=\scriptsize,\c,anchor=west] at (\s+0.15,0.42) {\lab};}
\foreach \s/\ic in {0/moon,2/cloud,4/sun}{\pic at (\s+0.6,0.64) {\ic};}
\end{tikzpicture}
\caption{Controle de ganho do sensor de luz segundo~\eqref{eq:nakarushton}. Cada curva cobre duas ou três ordens de brilho; ao passar para um fundo mais claro, $\sigma$ cresce e a curva se desloca ao longo do eixo logarítmico. Juntas, as curvas que se deslocam cobrem a faixa inteira.}
\label{fig:adaptation}
\end{figure}

O controle tem um preço, conhecido do capítulo~\ref{sec:code}: o sensor não informa o brilho, e sim o brilho \emph{em relação ao fundo}. Uma entrada constante deixa aos poucos de provocar resposta, e cada mudança provoca. As entradas da máquina falam, desde o início, de mudanças e não de níveis, e nisso está a mesma economia de disparos da proposição~\ref{prop:economy}~\cite{Barlow1961}. Daí também os sensores de ligar e de desligar do capítulo~\ref{sec:glutcomputer}~\cite{Schiller1992}: uns informam que ficou mais claro, outros que ficou mais escuro.

Os engenheiros repetiram esse truque nas \emph{câmeras de eventos}. Essa câmera não tira quadros ao ritmo de um relógio: cada pixel, por conta própria e sem ciclo comum, emite um evento ``mais claro'' ou ``mais escuro'' quando o logaritmo do brilho muda por uma fração dada~\cite{Lichtsteiner2008}. É exatamente o código de dois fios liga--desliga do capítulo~\ref{sec:glutcomputer}, realizado em silício, e ele dá à câmera uma faixa e uma rapidez inalcançáveis para um sensor de imagem comum~\cite{Gallego2022}.

\section{O olho: compressão até caber no fio}

O olho tem cerca de $10^8$ sensores de luz~\cite{Curcio1990}, e os neurônios de saída, que levam a imagem ao cérebro, são cerca de um milhão. Antes que o sinal saia do olho, ele é comprimido cerca de cem vezes. Os principais truques de compressão já nos são conhecidos. Cada neurônio de saída responde à diferença entre o brilho numa pequena mancha e o brilho em volta dela: é a subtração dos vizinhos por inibição, como no capítulo~\ref{sec:gaba}. Trechos uniformes da imagem quase não custam nada, enquanto bordas e mudanças são transmitidas. Os diferentes neurônios de saída são especializados: uns informam que clareou, outros que escureceu, outros ainda o movimento numa certa direção (fig.~\ref{fig:direction}); no camundongo foram contados mais de trinta desses canais de saída~\cite{Baden2016}.

Qual é a capacidade do cabo resultante? No porquinho-da-índia, a partir de registros dos neurônios de saída, mediram-se cerca de $875$ mil bits por segundo para $\sim10^5$ desses neurônios; recalculando para o milhão de neurônios de saída humanos, obtém-se da ordem de $10^7$ bits por segundo~\cite{Koch2006}, o mesmo que um velho cabo de rede de dez megabits. Com uma taxa média de alguns disparos por segundo por neurônio, isso dá alguns bits por disparo, como já saía no capítulo~\ref{sec:code}.

\section{O ouvido: um banco de filtros}

O ouvido resolve outro problema: o som precisa ser decomposto em frequências. Um engenheiro poria um banco de filtros passa-faixa, e o ouvido faz exatamente isso, mecanicamente. A vibração sonora percorre uma membrana elástica cujas propriedades mudam suavemente ao longo do seu comprimento, e cada trecho ressoa na sua própria frequência. Os sensores assentados num trecho respondem apenas à sua banda (fig.~\ref{fig:filterbank}). As frequências se distribuem ao longo do lugar de forma aproximadamente logarítmica: cada oitava recebe uma fração parecida dos sensores. O número do sensor que disparou é a frequência: o código de lugar do capítulo~\ref{sec:code}.

\begin{figure}[t]
\centering
\begin{tikzpicture}[>=Stealth,x=1.35cm,y=2.2cm]
\draw[->,gray!70!black] (-0.3,0) -- (7.6,0) node[above left,font=\small,black] {frequência, Hz};
\draw[->,gray!70!black] (-0.3,0) -- (-0.3,1.15) node[left,font=\small,black] {resposta};
\foreach \k/\f in {0/125,1/250,2/500,3/1000,4/2000,5/4000,6/8000}{
  \draw[gray!60!black] (\k,-0.02) -- (\k,0.02); \node[below,font=\scriptsize] at (\k,0) {\f};
  \draw[thick,blue!60!black] plot[domain={max(\k-1.2,-0.3)}:{\k+0.7},samples=60] (\x,{1/(1+((\x-\k)/(\x<\k?0.45:0.2))^4)});}
% a piano keyboard: seven white keys per octave, like the octaves on the axis
\foreach \i in {0,...,48}{\draw[gray!50!black,fill=white,line width=0.3pt] ({\i/7},-0.44) rectangle ({(\i+1)/7},-0.26);}
\foreach \o in {0,...,6}{\foreach \j in {0,1,3,4,5}{\fill[black] ({\o+(\j+1)/7-0.035},-0.35) rectangle ({\o+(\j+1)/7+0.035},-0.26);}}
\end{tikzpicture}
\caption{O ouvido como banco de filtros passa-faixa, esquema. Cada curva é a resposta dos sensores de um trecho a sons de frequências diferentes; as frequências centrais estão separadas por uma oitava, como teclas numa escala logarítmica. Os filtros reais têm uma encosta íngreme nas altas frequências e suave nas baixas. Embaixo, um teclado: nele, como no ouvido, cada oitava recebe o mesmo espaço.}
\label{fig:filterbank}
\end{figure}

Os ressonadores do ouvido não são passivos. Parte dos sensores sacode a própria membrana no ritmo do som e amplifica as vibrações fracas milhares de vezes em amplitude (em $66$--$76$\,dB), e com o aumento do volume o ganho cai~\cite{Ruggero1997,Robles2001}. Para o engenheiro, é um amplificador com realimentação positiva colocado bem perto do limite da autoexcitação: a mesma vida no limite da rede do capítulo~\ref{sec:gaba}; perto do limite, o sistema é especialmente sensível a sinais pequenos. A compressão de volume aqui é feita pelo mesmo amplificador: o que é baixo é amplificado muito, o que é alto, pouco.

O ouvido tem ainda um segundo código. Nas frequências baixas, os disparos dos neurônios vêm em fase com o som: o neurônio dispara num certo trecho de cada onda, embora não em toda onda. No porquinho-da-índia, o travamento de fase começa a enfraquecer em torno de $600$\,Hz e ainda é perceptível até cerca de $3.5$\,kHz~\cite{PalmerRussell1986}. Assim, os instantes dos disparos guardam o tempo exato do som e, a partir dele, a diferença de chegada do som aos dois ouvidos, pela qual a rede do capítulo~\ref{sec:glutcomputer} encontra a direção~\cite{Grothe2010}. Nas frequências altas os disparos não acompanham a onda, e resta apenas o código de lugar. Um mesmo ouvido usa os dois códigos do capítulo~\ref{sec:code}: o tempo onde os neurônios conseguem acompanhar e o lugar onde não conseguem.

\section{Os demais sensores}

\begin{table}[t]
\centering
\footnotesize
\setlength{\tabcolsep}{3pt}
\renewcommand{\arraystretch}{1.15}
\begin{tabular}{>{\raggedright}p{1.9cm}>{\raggedright}p{2.6cm}>{\raggedright}p{4.0cm}>{\raggedright\arraybackslash}p{4.6cm}}
\toprule
Sentido & O que mede & Como é a entrada & Truque do livro \\
\midrule
visão & luz & $\sim10^8$ sensores, compressão de $\sim100$ vezes, $\sim10^7$ bit/s & controle de ganho, subtração dos vizinhos, dois fios, direção do movimento \\
audição & pressão do ar & banco de filtros mecânico, amplificador ativo & código de lugar e código temporal, atrasos \\
tato & pressão, vibração & sensores de adaptação rápida e lenta & par ``filtro passa-altas --- filtro passa-baixas'' \\
temperatura & calor & sensores separados de calor e de frio & código de dois fios \\
olfato & moléculas & $\sim400$ tipos de receptores, um tipo por sensor & código combinatório: o cheiro é o conjunto de tipos que dispararam \\
paladar & substâncias no alimento & cinco qualidades: doce, amargo, salgado, azedo, umami & linhas rotuladas; parte das células libera ATP ou serotonina, e não glutamato \\
posição do corpo & estiramento dos músculos & sensores de comprimento e de velocidade de estiramento & realimentação para o regulador do movimento \\
\bottomrule
\end{tabular}
\caption{Como os sensores estão ligados. Todos os dispositivos da terceira coluna foram medidos; a coluna da direita os liga aos truques dos capítulos anteriores.}
\label{tab:sensors}
\end{table}

Os demais sentidos estão reunidos na tabela~\ref{tab:sensors}, e em quase toda linha se reconhece um truque dos capítulos anteriores. O tato usa um par de filtros: alguns sensores de pressão respondem ao toque e logo se calam, outros mantêm a resposta enquanto a pressão existe. A temperatura é transmitida por dois fios, separados para o calor e para o frio. A entrada mais incomum é o olfato. O ser humano tem cerca de quatrocentos tipos de receptores olfativos, e cada neurônio olfativo carrega um só tipo~\cite{Mombaerts2004}. Um cheiro não ativa um tipo, e sim um conjunto, e a mensagem é \emph{quais} tipos dispararam. Para o engenheiro, é um vetor binário de comprimento cerca de quatrocentos, com um número astronômico de valores possíveis.

\begin{keyblock}
\begin{proposition}[Entradas da máquina]
\label{prop:sensors}
Cada sentido é um transdutor que trabalha no limite físico de sensibilidade, comprime uma enorme faixa dinâmica por controle automático de ganho e transmite pelo barramento \nt{GLUT} sobretudo as \emph{mudanças}. Para isso, os sensores usam os mesmos truques da própria máquina: o código de dois fios, o código de lugar, o código temporal e a divisão pelo fundo. A construção dos sensores foi medida; que os seus códigos estejam ajustados para a máxima informação por disparo é uma hipótese bem fundamentada.
\end{proposition}
\end{keyblock}

As entradas, como todo o resto, funcionam de forma assíncrona. Cada sensor emite um disparo quando tem algo a informar, e os diferentes sentidos chegam com atrasos diferentes: o som em milissegundos, a luz em dezenas de milissegundos. Como a máquina os junta numa única imagem do mundo, sem um relógio comum, é um dos problemas de coordenação no tempo do capítulo~\ref{sec:synthesis}.

\section{Exercícios}

\begin{exercise}[\lvlA]
\label{ex:11:1}
\emph{Quanto tempo acumular fótons.} No crepúsculo chegam ao sensor $100$ fótons por segundo; um acréscimo é perceptível se for $3$ vezes maior que a flutuação~\eqref{eq:photonnoise}. (a)~Quanto tempo um sensor precisa para perceber uma mudança de brilho de $10\,\%$? (b)~Quantos sensores é preciso somar para conseguir em $0.2\,\text{s}$, e quantas vezes cai a resolução espacial?
\end{exercise}

\begin{exercise}[\lvlA]
\label{ex:11:2}
\emph{Quantos bits há num disparo.} (a)~O olho transmite $\sim10^7$ bit/s por $10^6$ neurônios de saída com taxa de $3$--$5$ disparos por segundo. Que fração da capacidade~\eqref{eq:spikeentropy} com $\Delta t=1\ms$ ele usa? (b)~Um cheiro ativa $20$ de $\approx400$ tipos de receptores. Quantos bits leva esse conjunto e quantos tipos em comum têm, em média, dois cheiros ao acaso?
\end{exercise}

\begin{exercise}[\lvlB]
\label{ex:11:3}
\emph{O que uma só curva~\eqref{eq:nakarushton} consegue.} (a)~Em que faixa de brilho a resposta sobe de $10\,\%$ a $90\,\%$ de $r_{\max}$? Quantas ordens de grandeza são? (b)~Quantas posições da curva deslocável são necessárias para cobrir dez ordens de brilho? E doze ordens de volume?
\end{exercise}

\begin{exercise}[\lvlB]
\label{ex:11:4}
\emph{O limite do código temporal no ouvido.} A diferença de chegada do som aos dois ouvidos vai até $0.22\,\text{m}/343\,\text{m/s}$. (a)~Os disparos vêm em fase com o tom: até que frequência a direção é determinada por eles sem ambiguidade? (b)~O disparo flutua $0.5\ms$, e é preciso distinguir $20$~\textmu s. Quantos neurônios de $100$ disparos por segundo é preciso promediar em $50\ms$?
\end{exercise}

\begin{exercise}[\lvlB]
\label{ex:11:5}
\emph{Uma câmera de eventos no cabo do olho.} Uma câmera de $10^6$ pixels com saída de $10^7$ bit/s envia um evento (endereço e sinal) quando $\ln I$ num pixel muda de $\ln1.2$. Com que velocidade máxima pode correr uma borda de $500$ pixels com salto de brilho $4$? Quantos quadros por segundo uma câmera comum, com $8$ bits por pixel, daria pela mesma saída?
\end{exercise}

\begin{exercise}[\lvlB]
\label{ex:11:6}
\emph{Simulação: controle de ganho.} O sensor~\eqref{eq:nakarushton} ajusta $\sigma$ com $\tau=1\,\text{s}$ em logaritmos, $\tau\,\dd\ln\sigma/\dd t=\ln I-\ln\sigma$, ou linearmente, $\tau\,\dd\sigma/\dd t=I-\sigma$; o fundo salta $1\to100\to1$. Após quantos segundos a resposta a um teste de $+10\,\%$ volta à metade da resposta adaptada, depois do salto para cima e para baixo, com cada lei?
\end{exercise}

\begin{exercise}[\lvlC]
\label{ex:11:7}
\emph{Para que mundo a curva está ajustada.} Com um ruído de saída pequeno e constante, a informação sobre o brilho é máxima quando $r(I)=r_{\max}\Phi_I(I)$, em que $\Phi_I$ é a função de distribuição dos brilhos. (a)~Para que distribuição a curva~\eqref{eq:nakarushton} é ótima, e qual a dispersão de $\log_{10}I$? (b)~Que curva é ótima com ruído de saída de Poisson?
\end{exercise}
